\documentclass{article} % For LaTeX2e
\usepackage{iclr2026_conference,times}

% Optional math commands from https://github.com/goodfeli/dlbook_notation.
\input{math_commands.tex}

\usepackage{hyperref}
\usepackage{url}
\usepackage{amsmath}
\usepackage{amssymb}
\usepackage{mathtools}
\usepackage{booktabs}
\usepackage{tabularx}
\usepackage{array}
\usepackage{algorithm}
\usepackage{algpseudocode}
\usepackage{graphicx}
\usepackage{multirow}
\usepackage{colortbl,xcolor}
\usepackage{makecell}
\definecolor{indomgray}{gray}{0.9}
\newcommand{\cellgray}{\cellcolor{indomgray}}
\newcolumntype{Y}{>{\centering\arraybackslash}X}
\title{Reinforced Flow Matching}

% Authors must not appear in the submitted version. They should be hidden
% as long as the \iclrfinalcopy macro remains commented out below.
% Non-anonymous submissions will be rejected without review.

\author{Anonymous Authors}

\newcommand{\fix}{\marginpar{FIX}}
\newcommand{\new}{\marginpar{NEW}}

%\iclrfinalcopy % Uncomment for camera-ready version, but NOT for submission.
\begin{document}

\maketitle

\begin{abstract}
The abstract.
\end{abstract}

\section{Introduction}
\label{sec:intro}

\section{Related work}
\label{sec:related_work}


\section{Method}

\subsection{Flow Matching}

Let $\mathbf{x}_0\sim p_{\mathcal{D}}$ denote a clean sample from the
data distribution and let
$\boldsymbol{\epsilon}\sim\mathcal{N}(\mathbf{0},\mathbf{I})$ denote
Gaussian noise. We consider the linear probability path
\begin{align}
    \mathbf{x}_t
    =
    (1-t)\mathbf{x}_0
    +
    t\boldsymbol{\epsilon},
    \qquad
    t\sim\mathcal{U}(0,1),
\end{align}
with the corresponding conditional velocity
\begin{align}
    \mathbf{v}_t
    (\mathbf{x}_t\mid\mathbf{x}_0,\boldsymbol{\epsilon})
    =
    \boldsymbol{\epsilon}-\mathbf{x}_0.
\end{align}

A velocity model $\mathbf{v}_{\boldsymbol{\phi}}$ is trained using the
flow matching objective
\begin{align}
\mathcal{L}_{\mathrm{FM}}(\boldsymbol{\phi})
=
\mathbb{E}_{t}
\mathbb{E}_{\mathbf{x}_0\sim p_{\mathcal{D}},
\boldsymbol{\epsilon}\sim\mathcal{N}(\mathbf{0},\mathbf{I})}
\left[
\left\|
\mathbf{v}_{\boldsymbol{\phi}}(\mathbf{x}_t,t)
-
\mathbf{v}_t
(\mathbf{x}_t\mid\mathbf{x}_0,\boldsymbol{\epsilon})
\right\|_2^2
\right].
\end{align}

Under the squared-error objective, the optimal marginal velocity field
is the conditional average of all velocities passing through
the same intermediate state:
\begin{align}
\mathbf{v}^{*}(\mathbf{x}_t,t)
=
\mathbb{E}_{\mathbf{x}_0,\boldsymbol{\epsilon}}
\left[
\mathbf{v}_t
(\mathbf{x}_t\mid\mathbf{x}_0,\boldsymbol{\epsilon})
\mid
\mathbf{x}_t,t
\right].
\label{eq:optimal-fm-velocity}
\end{align}

For an empirical dataset, this conditional expectation can be written
as
\begin{align}
\mathbf{v}^{*}(\mathbf{x}_t,t)
=
\sum_{\mathbf{x}_0\in\mathcal{D}}
p(\mathbf{x}_0\mid\mathbf{x}_t,t)
\,
\mathbb{E}_{\boldsymbol{\epsilon}}
\left[
\mathbf{v}_t
(\mathbf{x}_t\mid\mathbf{x}_0,\boldsymbol{\epsilon})
\mid
\mathbf{x}_0,\mathbf{x}_t,t
\right].
\label{eq:empirical-optimal-velocity}
\end{align}

Therefore, the optimal velocity is not a uniform average over all clean
samples. The contribution of each sample depends on its posterior
probability conditioned on the current intermediate state
$(\mathbf{x}_t,t)$.

\subsection{Reward-Reweighted Flow Matching}

Our goal is to adjust the probability mass assigned to a set of samples
according to their rewards. We assume that every sample initially has
unit weight. For a selected set $\mathcal{D}_r$, we change the weight of
each sample $\mathbf{x}_0\in\mathcal{D}_r$ from $1$ to
$W(\mathbf{x}_0)$, while samples outside $\mathcal{D}_r$ retain their
original unit weights.

The total weighted mass of $\mathcal{D}_r$ is
\begin{align}
S_w(\mathcal{D}_r)
=
S(\mathcal{D}_r)
\mathbb{E}_{\mathbf{x}_0\sim p_{\mathcal{D}_r}}
\left[
W(\mathbf{x}_0)
\right].
\label{eq:weighted-size}
\end{align}

The exact reweighted velocity requires the state-conditioned
contribution of every clean sample. To obtain a tractable expression,
we approximate the posterior proportions conditioned on
$(\mathbf{x}_t,t)$ using the global empirical proportions of
$\mathcal{D}\setminus\mathcal{D}_r$ and $\mathcal{D}_r$. Under this
dataset-level approximation, the reweighted velocity is
\begin{align}
\mathbf{v}_r^*(\mathbf{x}_t,t)
\approx~&
\frac{
S(\mathcal{D}\setminus\mathcal{D}_r)
}{
S(\mathcal{D}\setminus\mathcal{D}_r)
+
S_w(\mathcal{D}_r)
}
\mathbb{E}_{\mathbf{x}_0\sim
p_{\mathcal{D}\setminus\mathcal{D}_r}}
\mathbb{E}_{\boldsymbol{\epsilon}}
\left[
\mathbf{v}_t
(\mathbf{x}_t\mid\mathbf{x}_0,\boldsymbol{\epsilon})
\mid
\mathbf{x}_0,\mathbf{x}_t,t
\right]
\nonumber\\
&+
\frac{
S(\mathcal{D}_r)
}{
S(\mathcal{D}\setminus\mathcal{D}_r)
+
S_w(\mathcal{D}_r)
}
\mathbb{E}_{\mathbf{x}_0\sim p_{\mathcal{D}_r}}
\left[
W(\mathbf{x}_0)
\cdot
\mathbb{E}_{\boldsymbol{\epsilon}}
\left[
\mathbf{v}_t
(\mathbf{x}_t\mid\mathbf{x}_0,\boldsymbol{\epsilon})
\mid
\mathbf{x}_0,\mathbf{x}_t,t
\right]
\right].
\label{eq:initial-reweighted-velocity}
\end{align}

Equation~\eqref{eq:initial-reweighted-velocity} directly represents our
motivation. Samples outside $\mathcal{D}_r$ retain their original mass,
whereas every sample in $\mathcal{D}_r$ changes its mass from $1$ to
$W(\mathbf{x}_0)$.

Using
\begin{align}
&S(\mathcal{D}\setminus\mathcal{D}_r)
\mathbb{E}_{\mathbf{x}_0\sim
p_{\mathcal{D}\setminus\mathcal{D}_r}}
[\cdot]
\nonumber=
S(\mathcal{D})
\mathbb{E}_{\mathbf{x}_0\sim p_{\mathcal{D}}}
[\cdot]
-
S(\mathcal{D}_r)
\mathbb{E}_{\mathbf{x}_0\sim p_{\mathcal{D}_r}}
[\cdot],
\label{eq:set-decomposition}
\end{align}
we can rewrite
Eq.~\eqref{eq:initial-reweighted-velocity} as
\begin{align}
\mathbf{v}_r^*(\mathbf{x}_t,t)
\approx~&
\frac{
S(\mathcal{D})
}{
S(\mathcal{D})-S(\mathcal{D}_r)+S_w(\mathcal{D}_r)
}
\mathbb{E}_{\mathbf{x}_0\sim p_{\mathcal{D}}}
\mathbb{E}_{\boldsymbol{\epsilon}}
\left[
\mathbf{v}_t
(\mathbf{x}_t\mid\mathbf{x}_0,\boldsymbol{\epsilon})
\mid
\mathbf{x}_0,\mathbf{x}_t,t
\right]
\nonumber\\
&+
\frac{
S(\mathcal{D}_r)
}{
S(\mathcal{D})-S(\mathcal{D}_r)+S_w(\mathcal{D}_r)
}
\mathbb{E}_{\mathbf{x}_0\sim p_{\mathcal{D}_r}}
\left[
\left(
W(\mathbf{x}_0)-1
\right)
\cdot
\mathbb{E}_{\boldsymbol{\epsilon}}
\left[
\mathbf{v}_t
(\mathbf{x}_t\mid\mathbf{x}_0,\boldsymbol{\epsilon})
\mid
\mathbf{x}_0,\mathbf{x}_t,t
\right]
\right].
\label{eq:reweighted-full-decomposition}
\end{align}

We define
\begin{align}
\beta
&=
\frac{S(\mathcal{D}_r)}{S(\mathcal{D})},
\label{eq:beta-definition}
\\
\delta_r
&=
\mathbb{E}_{\mathbf{x}_0\sim p_{\mathcal{D}_r}}
\left[
W(\mathbf{x}_0)-1
\right].
\label{eq:delta-definition}
\end{align}

From Eq.~\eqref{eq:weighted-size}, we have
\begin{align}
S_w(\mathcal{D}_r)
=
S(\mathcal{D}_r)
\left(
1+\delta_r
\right).
\end{align}
Therefore,
\begin{align}
&S(\mathcal{D})-S(\mathcal{D}_r)+S_w(\mathcal{D}_r)=
S(\mathcal{D})(1+\beta\delta_r).
\label{eq:normalization-denominator}
\end{align}

Using the same dataset-level approximation to replace the full-data
average by the original optimal velocity, we obtain
\begin{align}
\mathbf{v}_r^*(\mathbf{x}_t,t)
\approx~&
\frac{1}{1+\beta\delta_r}
\mathbf{v}^{*}(\mathbf{x}_t,t)
\nonumber\\
&+
\frac{\beta}{1+\beta\delta_r}
\mathbb{E}_{\mathbf{x}_0\sim p_{\mathcal{D}_r}}
\left[
\left(
W(\mathbf{x}_0)-1
\right)
\cdot
\mathbb{E}_{\boldsymbol{\epsilon}}
\left[
\mathbf{v}_t
(\mathbf{x}_t\mid\mathbf{x}_0,\boldsymbol{\epsilon})
\mid
\mathbf{x}_0,\mathbf{x}_t,t
\right]
\right].
\label{eq:reweighted-velocity-decomposition}
\end{align}

We then use
\begin{align}
\frac{1}{1+\beta\delta_r}
=
1-
\frac{\beta\delta_r}{1+\beta\delta_r}
\end{align}
to write
\begin{align}
\mathbf{v}_r^*(\mathbf{x}_t,t)
\approx~&
\mathbf{v}^{*}(\mathbf{x}_t,t)
-
\frac{\beta\delta_r}{1+\beta\delta_r}
\mathbf{v}^{*}(\mathbf{x}_t,t)
\nonumber\\
&+
\frac{\beta}{1+\beta\delta_r}
\mathbb{E}_{\mathbf{x}_0\sim p_{\mathcal{D}_r}}
\left[
\left(
W(\mathbf{x}_0)-1
\right)
\cdot
\mathbb{E}_{\boldsymbol{\epsilon}}
\left[
\mathbf{v}_t
(\mathbf{x}_t\mid\mathbf{x}_0,\boldsymbol{\epsilon})
\mid
\mathbf{x}_0,\mathbf{x}_t,t
\right]
\right].
\label{eq:separated-normalization}
\end{align}

Since
\begin{align}
\delta_r\mathbf{v}^{*}(\mathbf{x}_t,t)
=
\mathbb{E}_{\mathbf{x}_0\sim p_{\mathcal{D}_r}}
\left[
\left(
W(\mathbf{x}_0)-1
\right)
\mathbf{v}^{*}(\mathbf{x}_t,t)
\right],
\end{align}
the latter two terms can be combined, yielding
\begin{align}
\mathbf{v}_r^*(\mathbf{x}_t,t)
\approx~&
\mathbf{v}^{*}(\mathbf{x}_t,t)
\nonumber\\
&+
\frac{\beta}{1+\beta\delta_r}
\mathbb{E}_{\mathbf{x}_0\sim p_{\mathcal{D}_r}}
\mathbb{E}_{\boldsymbol{\epsilon}}
\Big[
\left(
W(\mathbf{x}_0)-1
\right)
\cdot
\left(
\mathbf{v}_t
(\mathbf{x}_t\mid\mathbf{x}_0,\boldsymbol{\epsilon})
-
\mathbf{v}^{*}(\mathbf{x}_t,t)
\right)
\mid
\mathbf{x}_0,\mathbf{x}_t,t
\Big].
\label{eq:reward-reweighted-final}
\end{align}

Equation~\eqref{eq:reward-reweighted-final} shows that changing the
sample weight from $1$ to $W(\mathbf{x}_0)$ modifies the original
velocity through the corresponding sample-specific velocity
difference. When $W(\mathbf{x}_0)>1$, the sample receives additional
mass and its conditional velocity is reinforced. When
$W(\mathbf{x}_0)<1$, its original mass is reduced and the model moves
away from its conditional velocity. When $W(\mathbf{x}_0)=1$, the
sample produces no change.

\paragraph{Approximation assumptions.}

The derivation above contains several approximations. First, the exact
optimal velocity in Eq.~\eqref{eq:optimal-fm-velocity} is conditioned
on the current state $(\mathbf{x}_t,t)$. In
Eq.~\eqref{eq:initial-reweighted-velocity}, we replace the corresponding
state-dependent posterior proportions with global dataset proportions.
As a result, $\beta$ and $\delta_r$ are treated as independent of
$\mathbf{x}_t$ and $t$.


\subsection{Sample-Wise Reweighted Target}

For a clean rollout sample $\mathbf{x}_i^0$, we sample
$\boldsymbol{\epsilon}_i\sim\mathcal{N}(\mathbf{0},\mathbf{I})$ and
construct
\begin{align}
\mathbf{x}_i^{t_i}
=
(1-t_i)\mathbf{x}_i^0
+
t_i\boldsymbol{\epsilon}_i.
\end{align}

The old model predicts
\begin{align}
\mathbf{v}_i^{\mathrm{old}}
=
\mathbf{v}_{\boldsymbol{\theta}_{\mathrm{old}}}
(\mathbf{x}_i^{t_i},c_i,t_i),
\end{align}
from which its clean-sample prediction is
\begin{align}
\widehat{\mathbf{x}}_{0,i}^{\mathrm{old}}
=
\mathbf{x}_i^{t_i}
-
t_i\mathbf{v}_i^{\mathrm{old}}.
\end{align}

For the sampled linear path,
\begin{align}
\mathbf{x}_i^0
=
\mathbf{x}_i^{t_i}
-
t_i
\mathbf{v}_{t_i}
(\mathbf{x}_i^{t_i}
\mid
\mathbf{x}_i^0,\boldsymbol{\epsilon}_i).
\end{align}

Using one rollout sample to estimate the expectation in
Eq.~\eqref{eq:reward-reweighted-final}, the reweighted clean-sample
target becomes
\begin{align}
\overline{\mathbf{x}}_{0,i}
=
\widehat{\mathbf{x}}_{0,i}^{\mathrm{old}}
+
\frac{
\beta
\left(
W(\mathbf{x}_i^0,c_i)-1
\right)
}{
1+\beta\delta_r(c_i)
}
\left(
\mathbf{x}_i^0
-
\widehat{\mathbf{x}}_{0,i}^{\mathrm{old}}
\right).
\label{eq:reweighted-clean-target}
\end{align}

For a sample with $W(\mathbf{x}_i^0,c_i)>1$, the target moves from the
old prediction toward the clean rollout sample. For a sample with
$W(\mathbf{x}_i^0,c_i)<1$, the coefficient is negative and the target
moves away from the rollout sample.

\subsection{Reward-Based Sample Weighting}

For the $K$ rollout samples generated under the same prompt $c$, we
first standardize and clip their raw rewards:
\begin{align}
\overline{R}_c
&=
\frac{1}{K}
\sum_{k=1}^{K}R_k,
\\
Z_c
&=
\sqrt{
\frac{1}{K}
\sum_{k=1}^{K}
(R_k-\overline{R}_c)^2
+
\epsilon
},
\\
\widehat{A}_k
&=
\frac{1}{A}
\operatorname{clip}
\left(
\frac{R_k-\overline{R}_c}{Z_c},
-A,A
\right).
\label{eq:normalized-reward}
\end{align}
This normalization ensures that
$\widehat{A}_k\in[-1,1]$, where positive and negative values indicate
above- and below-average rewards, respectively. We consider the
following two alternative mappings from the normalized reward to the
sample weight.

\paragraph{Exponential weighting.}
The first option uses an exponential mapping:
\begin{align}
W_{\mathrm{exp}}(\mathbf{x}_0^k,c)
=
\exp(\widehat{A}_k).
\label{eq:exponential-reward-weight}
\end{align}
A sample with $\widehat{A}_k=0$ retains its original unit weight,
while samples with positive and negative normalized rewards receive
weights above and below one, respectively. Since
$\widehat{A}_k\in[-1,1]$, the resulting weights satisfy
\begin{align}
W_{\mathrm{exp}}(\mathbf{x}_0^k,c)
\in
[e^{-1},e].
\end{align}
This mapping guarantees strictly positive weights and provides a
multiplicative adjustment of the original sample mass.

\paragraph{Linear weighting.}
Alternatively, we directly add the normalized reward to the original
unit weight:
\begin{align}
W_{\mathrm{lin}}(\mathbf{x}_0^k,c)
=
1+\widehat{A}_k.
\label{eq:linear-reward-weight}
\end{align}
The resulting weights satisfy
\begin{align}
W_{\mathrm{lin}}(\mathbf{x}_0^k,c)
\in
[0,2].
\end{align}
This mapping provides a direct additive interpretation:
$\widehat{A}_k=0$ preserves the original weight,
$\widehat{A}_k=1$ doubles the sample weight, and
$\widehat{A}_k=-1$ reduces the sample weight to zero. In practice, if
strictly positive weights are required for numerical stability, the
linear mapping can be implemented as
\begin{align}
W_{\mathrm{lin}}(\mathbf{x}_0^k,c)
=
\max\left\{
\epsilon,\,
1+\widehat{A}_k
\right\}.
\end{align}






\begin{algorithm}[H]
\caption{REWARD-REWEIGHTED FLOW MATCHING}
\label{alg:velocity-add-delete}
\begin{algorithmic}[1]
\Require Initial model $\theta_0$, reward function $R$, prompt distribution
$p(c)$, rollout size $K$, number of timesteps per sample $J$, batch size $B$,
reinforcement coefficient $\beta$, advantage clip $A$, reference regularization
$\lambda_{\mathrm{ref}}$, EMA rate $\eta$, and numerical constant $\epsilon$.

\State $\theta,\theta_{\mathrm{old}},\theta_{\mathrm{ref}}
\leftarrow\theta_0$.

\For{each training iteration}
    \State $\mathcal{D}\leftarrow\varnothing$.

    \For{each prompt $c\sim p(c)$}
        \State Sample $K$ clean rollout samples
        $\{\mathbf{x}_0^k\}_{k=1}^{K}$
        with $\theta_{\mathrm{old}}$ conditioned on $c$.

        \State Compute rewards
        $R_k\leftarrow R(\mathbf{x}_0^k,c)$
        for $k=1,\ldots,K$.

        \State $\overline{R}_c
        \leftarrow
        K^{-1}\sum_{k=1}^{K}R_k$.

        \State $Z_c
        \leftarrow
        \sqrt{
        K^{-1}\sum_{k=1}^{K}
        (R_k-\overline{R}_c)^2+\epsilon
        }$.

        \For{$k=1,\ldots,K$}
            \State $\widehat{A}_k
            \leftarrow
            \operatorname{clip}
            \left(
            \frac{R_k-\overline{R}_c}{Z_c},
            -A,A
            \right)/A$.
        
            \State $W(\mathbf{x}_0^k,c)
            \leftarrow
            \max\left\{\epsilon,\,1+\widehat{A}_k\right\}$.
        \EndFor

        \State $\displaystyle
        \delta_r(c)
        \leftarrow
        \frac{1}{K}
        \sum_{k=1}^{K}
        \left(
        W(\mathbf{x}_0^k,c)-1
        \right)$.

        % \For{$k=1,\ldots,K$}
        %     \State Sample or store $J$ training timesteps
        %     $\{t_{k,j}\}_{j=1}^{J}$ for sample $\mathbf{x}_0^k$.

        %     \For{$j=1,\ldots,J$}
        %         \State Add
        %         $\left\{
        %         c,\mathbf{x}_0^k,t_{k,j},
        %         W(\mathbf{x}_0^k,c),\delta_r(c)
        %         \right\}$
        %         to $\mathcal{D}$.
        %     \EndFor
        % \EndFor
    \EndFor

    \State Set the accumulated gradients of $\theta$ to zero.

    \For{each minibatch
    $\{
    (c_i,\mathbf{x}_i^0,t_i,
    W(\mathbf{x}_i^0,c_i),\delta_r(c_i))
    \}_{i=1}^{B}\sim\mathcal{D}$}

        \State Sample
        $\boldsymbol{\epsilon}_i\sim\mathcal{N}(\mathbf{0},\mathbf{I})$
        and set
        $\mathbf{x}_i^{t_i}
        \leftarrow
        (1-t_i)\mathbf{x}_i^0
        +
        t_i\boldsymbol{\epsilon}_i$
        for $i=1,\ldots,B$.

        \State $\mathbf{v}_i^{\theta}
        \leftarrow
        \mathbf{v}_{\theta}
        (\mathbf{x}_i^{t_i},c_i,t_i)$.

        \State $\mathbf{v}_i^{\mathrm{old}}
        \leftarrow
        \mathbf{v}_{\theta_{\mathrm{old}}}
        (\mathbf{x}_i^{t_i},c_i,t_i)$.

        \State $\mathbf{v}_i^{\mathrm{ref}}
        \leftarrow
        \mathbf{v}_{\theta_{\mathrm{ref}}}
        (\mathbf{x}_i^{t_i},c_i,t_i)$.

        \State $\widehat{\mathbf{x}}_{0,i}^{\theta}
        \leftarrow
        \mathbf{x}_i^{t_i}
        -
        t_i\mathbf{v}_i^{\theta}$.

        \State $\widehat{\mathbf{x}}_{0,i}^{\mathrm{old}}
        \leftarrow
        \mathbf{x}_i^{t_i}
        -
        t_i\mathbf{v}_i^{\mathrm{old}}$.

        \State $\displaystyle
        \overline{\mathbf{x}}_{0,i}
        \leftarrow
        \widehat{\mathbf{x}}_{0,i}^{\mathrm{old}}
        +
        \frac{
        \beta
        \left(
        W(\mathbf{x}_i^0,c_i)-1
        \right)
        }{
        1+\beta\delta_r(c_i)
        }
        \left(
        \mathbf{x}_i^0
        -
        \widehat{\mathbf{x}}_{0,i}^{\mathrm{old}}
        \right)$.

        \State $\displaystyle
        \mathcal{L}_{\mathrm{target}}
        \leftarrow
        \frac{A}{B}
        \sum_{i=1}^{B}
        \frac{
        \left\|
        \widehat{\mathbf{x}}_{0,i}^{\theta}
        -
        \overline{\mathbf{x}}_{0,i}
        \right\|_2^2
        }{
        \operatorname{sg}
        \left[
        \operatorname{mean}
        \left|
        \widehat{\mathbf{x}}_{0,i}^{\theta}
        -
        \overline{\mathbf{x}}_{0,i}
        \right|
        \right]
        }$.

        \State $\displaystyle
        \mathcal{L}_{\mathrm{ref}}
        \leftarrow
        \frac{1}{B}
        \sum_{i=1}^{B}
        \left\|
        \mathbf{v}_i^{\theta}
        -
        \mathbf{v}_i^{\mathrm{ref}}
        \right\|_2^2$.

        \State $\mathcal{L}
        \leftarrow
        \frac{1}{N_{acc}}(\mathcal{L}_{\mathrm{target}}
        +
        \lambda_{\mathrm{ref}}
        \mathcal{L}_{\mathrm{ref}})$.

        \State Backpropagate $\mathcal{L}$ and accumulate the gradients
        without updating $\theta$.
    \EndFor

    \State Update $\theta$ once using the accumulated gradients.

    \State $\theta_{\mathrm{old}}
    \leftarrow
    \eta\theta_{\mathrm{old}}
    +(1-\eta)\theta$.
\EndFor
\end{algorithmic}
\end{algorithm}



\newpage            

\section{Experiments}
\label{sec:experiments}

\subsection{Experimental setup}
\label{sec: exp-setting}
We follow the experimental setup used by DiffusionNFT~\citep{zheng2026diffusionnft}. Unless otherwise specified, all experiments start from SD3.5-Medium~\citep{esser2024scaling} at $512 \times 512$ resolution.

\paragraph{Reward models.}
We evaluate both rule-based and model-based reward optimization. The rule-based tasks are GenEval for compositional text-to-image alignment~\citep{ghosh2023geneval} and OCR visual text rendering, using the same prompt splits and partial-credit reward design as the DiffusionNFT and FlowGRPO. The model-based rewards include PickScore~\citep{kirstain2023pickapic}, CLIPScore~\citep{hessel2021clipscore}, HPSv2.1~\citep{wu2023human}, Aesthetic score~\citep{schuhmann2022laion}, ImageReward~\citep{xu2023imagereward}, and UnifiedReward~\citep{wang2025unifiedreward}. Following DiffusionNFT, Pick-a-Pic prompts are used for preference-oriented training~\citep{kirstain2023pickapic}, while DrawBench prompts are reserved for out-of-domain evaluation~\citep{saharia2022photorealistic}.

\paragraph{Training protocol.}
We fine-tune the base model with LoRA~\citep{hu2022lora} using rank $r=32$, scaling factor $\alpha=64$, and learning rate $3 \times 10^{-4}$. Each training epoch contains 48 prompt groups, and each group samples $G=24$ images from the current rollout policy. Rewards are normalized within each group before constructing the optimization weights, so methods are compared under the same reward scale. For single-reward comparisons and ablations, rollouts use 10 sampling steps to match the efficient setting in DiffusionNFT. For multi-reward training and final evaluation, we use 40 sampling steps to reduce sampler-induced quality differences.

\paragraph{Baselines and evaluation.}

\subsection{Main results}
\begin{table*}[t]
\centering
\caption{\textbf{Evaluation Results across tasks and reward settings.}
\colorbox{indomgray}{Gray-colored} cells indicate in-domain reward performance for each task.
$\dagger$ indicates evaluated results on official checkpoints.
$\ddagger$ indicates evaluated under $1024 \times 1024$ resolution. \textbf{Bold} indicates the best result within each task block. For models in Baselines category, we evaluate GenEval and OCR on its corresponding dataset, while other rewards are evaluated on DrawBench \citep{saharia2022photorealistic}. References for models: FlowGRPO \citep{liu2025flowgrpo}, AWM \citep{xue2026advantage}, DiffusionNFT \citep{zheng2026diffusionnft}, Rethinking \citep{choi2026rethinking}.}
\label{tab:eval_results}
\setlength{\tabcolsep}{6pt}
\renewcommand{\arraystretch}{0.95}
\footnotesize
\resizebox{\textwidth}{!}{
\begin{tabular}{ll ccccccc}
\toprule
Eval. Dataset & Model & GenEval & OCR & PickScore & ClipScore & HPSv2.1 & Aesthetic & ImgRwd \\
\midrule
\multirow{5}{*}{Baselines} 
& SD-XL$^\ddagger$     & 0.55 & 0.14 & 22.42 & 0.287 & 0.280 & 5.60 & 0.76 \\
& SD3.5-L$^\ddagger$   & 0.71 & 0.68 & 22.91 & 0.289 & 0.288 & 5.50 & 0.96 \\
& FLUX.1-Dev          & 0.66 & 0.59 & 22.84 & 0.295 & 0.274 & 5.71 & 0.96 \\
& SD3.5-M             & 0.24 & 0.12 & 20.51 & 0.237 & 0.204 & 5.13 & -0.58 \\
& \quad + CFG          & 0.63 & 0.59 & 22.34 & 0.285 & 0.279 & 5.36 & 0.85 \\
\midrule
\multirow{5}{*}{GenEval} 
& FlowGRPO$^\dagger$   & \cellgray 0.95 & - & 22.51 & 0.293 & 0.274 & 5.32 & 1.06 \\
& AWM                 & \cellgray 0.89 & - & 22.00 & 0.302 & 0.242 & 4.94 & 0.84 \\
& DiffusionNFT        & \cellgray 0.95 & - & 22.88 & 0.303 & 0.289 & 5.25 & 1.21 \\
& Rethinking (ELBO-ODE) & \cellgray 0.96 & - & 22.85 & 0.305 & 0.289 & 5.33 & 1.26 \\
& (Ours)              & \cellgray - & - & - & - & - & - & - \\
\midrule
\multirow{5}{*}{OCR} 
& FlowGRPO$^\dagger$   & - & \cellgray 0.92 & 22.41 & 0.290 & 0.280 & 5.32 & 0.95 \\
& AWM                 & - & \cellgray 0.80 & 20.70 & 0.301 & 0.206 & 4.53 & -0.13 \\
& DiffusionNFT        & - & \cellgray 0.93 & \cellgray 22.09 & \cellgray 0.307 & \cellgray 0.277 & 5.17 & 0.97 \\
& Rethinking (ELBO-ODE) & - & \cellgray 0.94 & \cellgray 22.93 & \cellgray 0.315 & \cellgray 0.302 & 5.33 & 1.34 \\
& (Ours) & - & \cellgray - & \cellgray - & \cellgray - & \cellgray - & - & - \\
\midrule
\multirow{4}{*}{Drawbench} 
& FlowGRPO$^\dagger$   & - & - & \cellgray 23.50 & 0.280 & 0.316 & 5.90 & 1.29 \\
& DiffusionNFT        & - & - & \cellgray 23.61 & \cellgray 0.288 & \cellgray 0.344 & 6.04 & 1.46 \\
& Rethinking (ELBO-ODE) & - & - & \cellgray 23.68 & \cellgray 0.296 & \cellgray 0.325 & 6.06 & 1.45 \\
& (Ours) & - & - & \cellgray - & \cellgray - & \cellgray - & - & - \\
\bottomrule
\end{tabular}
}
\end{table*}

\subsection{Ablation study}




We conduct ablations to understand the contribution and robustness of our key design choices. Specifically, we study the add-delete target, adaptive loss weighting, rollout-policy update, and target-editing strength under GenEval single-reward training. 
Unless stated otherwise, every variant uses the same SD3.5-Medium, prompts, LoRA configuration, group-normalized rewards, rollout
budget, and evaluation sampler described in Section~\ref{sec: exp-setting}. 



\paragraph{Reweighting coefficient.}
We ablate the reinforcement coefficient $\beta$, which controls the magnitude
of the reweighting. We keep all other training
settings fixed, and compare $\beta\in\{0.1,1,10\}$. 

\begin{table}[H]
\centering
\caption{Ablation of the reinforcement coefficient $\beta$ on GenEval.}
\label{tab:ablation_beta}
\scriptsize
\setlength{\tabcolsep}{3pt}
\begin{tabularx}{0.5\columnwidth}{@{}X*{6}{c}@{}}
\toprule
Reinforcement coefficient
& \multicolumn{6}{c}{GenEval $\uparrow$ at training iteration} \\
\cmidrule(l){2-7}
& 0 & 100 & 200 & 300 & 400 & 500 \\
\midrule
$\beta=0.1$ & -- & -- & -- & -- & -- & -- \\
$\beta=1$   & -- & -- & -- & -- & -- & -- \\
$\beta=10$  & -- & -- & -- & -- & -- & -- \\
\bottomrule
\end{tabularx}
\end{table}





\paragraph{Reward-to-weight mapping.}
We ablate the choice of reward-to-weight mapping
$W(\mathbf{x}_0^k,c)$, which determines how the normalized reward
$\widehat{A}_k$ adjusts the original unit weight of each rollout sample.
We keep all other training settings fixed and compare the exponential
mapping in Eq.~\eqref{eq:exponential-reward-weight} with the linear
mapping in Eq.~\eqref{eq:linear-reward-weight}.

\begin{table}[H]
\centering
\caption{Ablation of the reward-to-weight mapping
$W(\mathbf{x}_0^k,c)$ on GenEval.}
\label{tab:ablation_weight_mapping}
\scriptsize
\setlength{\tabcolsep}{3pt}
\begin{tabularx}{\columnwidth}{@{}X*{6}{c}@{}}
\toprule
Reward-to-weight mapping
& \multicolumn{6}{c}{GenEval $\uparrow$ at training iteration} \\
\cmidrule(l){2-7}
& 0 & 100 & 200 & 300 & 400 & 500 \\
\midrule
Exponential weighting & -- & -- & -- & -- & -- & -- \\
Linear weighting      & -- & -- & -- & -- & -- & -- \\
\bottomrule
\end{tabularx}
\end{table}



\paragraph{Target components.} We first verify that the single target objective reproduces the original two directions in  DiffusionNFT. We then compare the full add-delete target
with addition-only and deletion-only variants. The latter three variants use the same signed reward decomposition and adaptive loss normalization, so their differences isolate the contribution of each target direction.

\begin{table}[H]
\centering
\caption{Ablation of the velocity target components on GenEval.
{\color{blue}(After the results are complete, convert to a line plot following DiffusionNFT.)}}
\label{tab:ablation_target}
\scriptsize
\setlength{\tabcolsep}{3pt}
\begin{tabularx}{\columnwidth}{@{}X*{6}{c}@{}}
\toprule
Variant & \multicolumn{6}{c}{GenEval $\uparrow$ at training iteration} \\
\cmidrule(l){2-7}
& 0 & 100 & 200 & 300 & 400 & 500 \\
\midrule
DiffusionNFT & -- & -- & -- & -- & -- & -- \\
Equivalent single target objective & -- & -- & -- & -- & -- & -- \\
\midrule
Addition only & -- & -- & -- & -- & -- & -- \\
Deletion only & -- & -- & -- & -- & -- & -- \\
% Full add-delete & -- & -- & -- & -- & -- & -- \\
\bottomrule
\end{tabularx}
\end{table}



\paragraph{Adaptive loss weighting.}
We keep the full add-delete target fixed and compare the self-normalized
$\hat{\mathbf{x}}_0$ regression with three fixed timestep-weighting functions.
This comparison measures whether adaptive normalization improves convergence
without changing the signed target displacement.

\begin{table}[H]
\centering
\caption{Ablation of loss weighting on GenEval.}
\label{tab:ablation_weighting}
\scriptsize
\setlength{\tabcolsep}{3pt}
\begin{tabularx}{\columnwidth}{@{}X*{6}{c}@{}}
\toprule
Loss weighting & \multicolumn{6}{c}{GenEval $\uparrow$ at training iteration} \\
\cmidrule(l){2-7}
& 0 & 100 & 200 & 300 & 400 & 500 \\
\midrule
$w(t)=1-t$ & -- & -- & -- & -- & -- & -- \\
$w(t)=1$ & -- & -- & -- & -- & -- & -- \\
$w(t)=t$ & -- & -- & -- & -- & -- & -- \\
Adaptive normalization & -- & -- & -- & -- & -- & -- \\
\bottomrule
\end{tabularx}
\end{table}





\paragraph{Decoupling soft updates for rollout and target construction.}
The EMA model in our full method serves two distinct roles: it generates
rollout samples for reward evaluation and provides the base velocity used to
construct the reward-reweighted training target. To identify where temporal
smoothing is most beneficial, we decouple these two roles and compare four
configurations. All
settings use the same reward weighting, reference regularization, number of
generated images, and optimization budget.

\begin{table}[H]
\centering
\caption{Ablation of soft updates for rollout generation and target
construction on GenEval. ``Latest'' denotes a detached snapshot of the model
parameters at the beginning of each training iteration, while ``EMA'' denotes
the soft-updated model.}
\label{tab:ablation_soft_update_roles}
\scriptsize
\setlength{\tabcolsep}{3pt}
\begin{tabularx}{\columnwidth}{@{}XX*{6}{c}@{}}
\toprule
Rollout policy & Target policy
& \multicolumn{6}{c}{GenEval $\uparrow$ at training iteration} \\
\cmidrule(l){3-8}
& & 0 & 100 & 200 & 300 & 400 & 500 \\
\midrule
EMA & EMA
& -- & -- & -- & -- & -- & -- \\

Latest & EMA
& -- & -- & -- & -- & -- & -- \\

EMA & Latest
& -- & -- & -- & -- & -- & -- \\

Latest & Latest
& -- & -- & -- & -- & -- & -- \\
\bottomrule
\end{tabularx}
\end{table}





\paragraph{Soft update of the rollout policy.}
We isolate the EMA update by comparing a fully on-policy update ($\eta=0$), a
fixed slow update ($\eta=0.9$), and the increasing schedule used by our full
configuration. All three settings use the full add-delete target and the same
number of generated images.

\begin{table}[H]
\centering
\caption{Ablation of rollout-policy updates on GenEval.
(After the results are complete, convert this table into a line plot following
DiffusionNFT.)}
\label{tab:ablation_ema}
\scriptsize
\setlength{\tabcolsep}{3pt}
\begin{tabularx}{\columnwidth}{@{}X*{6}{c}@{}}
\toprule
Update schedule & \multicolumn{6}{c}{GenEval $\uparrow$ at training iteration} \\
\cmidrule(l){2-7}
& 0 & 100 & 200 & 300 & 400 & 500 \\
\midrule
On-policy, $\eta=0$ & -- & -- & -- & -- & -- & -- \\
Fixed EMA, $\eta=0.9$ & -- & -- & -- & -- & -- & -- \\
Increasing EMA schedule & -- & -- & -- & -- & -- & -- \\
\bottomrule
\end{tabularx}
\end{table}

% \paragraph{Addition--deletion strength.}
% Finally, we first vary the symmetric strength $\beta_a=\beta_d$. We then fix
% $\beta_a+\beta_d=2$ and vary the relative strength of addition and deletion.
% The first comparison measures the trade-off between convergence speed and
% stability, while the second tests whether the two editing directions require
% comparable strength.

% \begin{table}[H]
% \centering
% \caption{Sensitivity to addition and deletion strengths on GenEval.
% (After the results are complete, convert this table into a line plot following
% DiffusionNFT.)}
% \label{tab:ablation_strength}
% \scriptsize
% \setlength{\tabcolsep}{2.5pt}
% \begin{tabularx}{\columnwidth}{@{}Xcc*{6}{c}@{}}
% \toprule
% Setting & $\beta_a$ & $\beta_d$ &
% \multicolumn{6}{c}{GenEval $\uparrow$ at training iteration} \\
% \cmidrule(l){4-9}
% & & & 0 & 100 & 200 & 300 & 400 & 500 \\
% \midrule
% \multirow{4}{*}{Symmetric}
% & $0.1$ & $0.1$ & -- & -- & -- & -- & -- & -- \\
% & $0.5$ & $0.5$ & -- & -- & -- & -- & -- & -- \\
% & $1.0$ & $1.0$ & -- & -- & -- & -- & -- & -- \\
% & $2.0$ & $2.0$ & -- & -- & -- & -- & -- & -- \\
% \midrule
% \multirow{3}{*}{Fixed sum}
% & $2/3$ & $4/3$ & -- & -- & -- & -- & -- & -- \\
% & $1.0$ & $1.0$ & -- & -- & -- & -- & -- & -- \\
% & $4/3$ & $2/3$ & -- & -- & -- & -- & -- & -- \\
% \bottomrule
% \end{tabularx}
% \end{table}




\section{Conclusion}
\label{sec:conclusion}

\bibliography{iclr2026_conference}
\bibliographystyle{iclr2026_conference}

\appendix
\section{Appendix}
You may include other additional sections here.

\end{document}
